\subsection{Elliptic representations separate interior points from cusps}

\begin{result}{E.192}{Proved}
{An actual pro-three elliptic detector over every finite field of characteristic different from three}
Let \(k=\mathbb F_q\) have characteristic different from three and put
\(U=\mathbb P^1_k\setminus\{0,1,\infty\}\).
For every \(a\in U(k)\) and every rational cusp \(c\), the section of
\(\pi_1(U)^{(3)}\to G_k\) defined by \(a\) is not geometrically conjugate
to a section whose image is contained in a decomposition group of \(c\).
Here the superscript means the actual geometrically pro-three quotient.
The assertion includes every tangential section and every choice of its
inertial parameter.

There is an explicit finite detector.  For an integer \(N\) such that
\[
                         3^N>2q+1,
\tag{E.192.1}
\]
one of the three elliptic representations below, reduced modulo \(3^N\),
separates the two section classes.  Its geometric image is a finite
three-group of order at most \(3^{3N-2}\).
Thus separation is already visible in an actual finite geometric
nilpotent quotient; it is not a consequence of replacing that quotient
by a rational unipotent completion.
Consequently the unique rational Frobenius path of E.135 has a first
nonintegral weight at most \(3N-2\).  If \(q\equiv-1\pmod3\), that
weight is even and is at most \(2\lfloor(3N-2)/2\rfloor\).
\end{result}

\begin{proof}
It suffices first to treat \(c=0\).  Over \(U\) consider the smooth
elliptic scheme
\[
                  \mathcal E_t:\quad y^2+3xy+t y=x^3 .
\tag{E.192.2}
\]
The integral Weierstrass formulas give
\[
 \begin{gathered}
 b_2=9,\quad b_4=3t,\quad b_6=t^2,\quad b_8=0,\\
 c_4=81-72t,\qquad
 \Delta=27t^3(1-t).
 \end{gathered}
\tag{E.192.3}
\]
These formulas also apply in characteristic two.  In particular the
fibres over \(U\) are smooth.  The section \(P=(0,0)\) has exact order
three: its tangent is \(y=0\), and this line has intersection divisor
\(3(P)\) on the smooth cubic.  It is not the origin at infinity.

The finite etale group schemes \(\mathcal E[3^n]\) give the actual
representation
\[
       r:\pi_1(U)\longrightarrow
                    \operatorname{GL}_2(\mathbb Z_3).
\tag{E.192.4}
\]
Its determinant is the cyclotomic character, by the Galois-equivariant
Weil pairing
\citep[III.8, Proposition 8.3, pp.97--98]{AuditSilverman2009}.
On the geometric subgroup that determinant is one.  Reduction modulo
three fixes the nonzero vector supplied by \(P\).  In a basis whose
first vector is \(P\), its geometric reduction is therefore contained in
\[
       \left\{\begin{pmatrix}1&b\\0&1\end{pmatrix}:
                      b\in\mathbb F_3\right\}.
\tag{E.192.5}
\]
This group has order three.  The kernel of reduction
\(\operatorname{SL}_2(\mathbb Z_3)\to
\operatorname{SL}_2(\mathbb F_3)\) is pro-three, as follows from its
successive congruence quotients.  Hence the entire geometric image of
\(r\) is pro-three.  The restriction of \(r\) to the geometric
fundamental group consequently kills the kernel of its maximal
pro-three quotient.  This proves that the \emph{actual arithmetic}
representation E.192.4 factors through \(\pi_1(U)^{(3)}\); an assertion
about an outer representation alone would not suffice here.

We next compute its restriction at the named cusp.  Over
\(K=k((t))\), the displayed integral model has unit \(c_4\),
discriminant of valuation three, and special fibre
\[
                         y(y+3x)=x^3 .
\tag{E.192.6}
\]
Its node has the two distinct rational tangent lines \(y=0\) and
\(y=-3x\).  Thus the reduction is split multiplicative, including in
characteristic two.  Tate uniformization gives a parameter
\(Q\in K^\times\), \(v_t(Q)>0\), and a Galois-equivariant
isomorphism \(\overline K^\times/Q^{\mathbb Z}
\simeq\mathcal E_t(\overline K)\)
\citep[C.14, Theorem 14.1(a)--(b), pp.444--445]{AuditSilverman2009}.
Taking compatible three-power torsion gives the exact sequence
\[
 0\longrightarrow\mathbb Z_3(1)
   \longrightarrow T_3\mathcal E_t
   \longrightarrow\mathbb Z_3\longrightarrow0.
\tag{E.192.7}
\]
The left term comes from roots of unity and the right term from the
exponents of compatible roots of \(Q\).  This proves the sequence
directly from the stated uniformization.

Restriction to the cusp decomposition group is precisely restriction
of the finite torsion covers to this punctured formal neighbourhood.
The representation is tame there: its inertial extension character is
the three-primary Kummer character of \(Q\), so wild inertia of
characteristic different from three acts trivially.  Consequently
E.192.7 descends to the same cusp decomposition group in the
geometrically pro-three fundamental group.  This is the required
local comparison, with no identification of unrelated path torsors.
Every element of this decomposition group above arithmetic Frobenius
has, in the Tate basis, a matrix
\[
                          \begin{pmatrix}q&b\\0&1\end{pmatrix}
                         \quad(b\in\mathbb Z_3).
\tag{E.192.8}
\]
In particular its trace is \(q+1\) and its determinant is \(q\).
Changing the tangential section or multiplying Frobenius by inertia
changes only the upper-right entry.

At an interior point \(a\), functoriality of the very same torsion
covers identifies the specialization with \(T_3(\mathcal E_a)\).
Write \(m_a=\#\mathcal E_a(k)\).  Frobenius on that Tate module has
\[
                 \det(F_a)=q,\qquad
                 \operatorname{tr}(F_a)=q+1-m_a .
\tag{E.192.9}
\]
For completeness, the determinant of an elliptic endomorphism on a
prime-to-characteristic Tate module is its degree.  The endomorphism
\(1-F_a\) is separable with kernel \(\mathcal E_a(k)\), so
\(\det(1-F_a)=m_a\); the two-by-two determinant identity gives
E.192.9.  These are
\citet[III.8, Proposition 8.6, p.99; V.2, Proposition 2.3 and
Theorem 2.3.1, pp.141--143]{AuditSilverman2009}.
The integer \(m_a\) is positive because the origin is rational.
Thus E.192.8 and E.192.9 have different traces in \(\mathbb Z_3\).
This already proves separation.  Equivalently, the good fibre has
\(T_3(\mathcal E_a)^{G_k}=0\), whereas the triangular cusp matrix has
a nonzero eigenvector of eigenvalue one.

There is no need to invoke a Weil bound for the finite assertion.
For each \(x\in k\), the Weierstrass equation has at most two values
of \(y\), and there is one point at infinity.  Hence
\(1\le m_a\le2q+1\).  Under E.192.1 its reduction modulo \(3^N\)
is nonzero, so the traces already differ there.  The geometric
modulo-\(3^N\) image \(H_N\) is contained in the inverse image of
E.192.5 inside \(\operatorname{SL}_2(\mathbb Z/3^N)\), which has order
\(3^{1+3(N-1)}\).  This proves the finite size bound.

If a characteristic geometric quotient is desired, put \(d=3N-2\)
and let \(\Pi\) be the free pro-three geometric group.  The finite
characteristic quotient
\[
     \Pi\big/\overline{\Pi^{3^d}\Gamma_{d+1}(\Pi)}
\tag{E.192.10}
\]
dominates \(H_N\): a group of order at most \(3^d\) has exponent
dividing \(3^d\) and nilpotency class at most \(d\).  Finiteness of
E.192.10 follows by successive lower central quotients of the
finitely generated nilpotent group of bounded exponent.  Its kernel
is characteristic, so the actual arithmetic extension descends as
well.  This supplies an explicit characteristic finite nilpotent
quotient if one does not wish to retain the specific elliptic image.
An integral conjugator in nilpotency class \(d\) would map to a
conjugator in \(H_N\), contradicting the trace calculation.  Thus the
unique rational path in E.135 fails integrality by weight \(d\).
Its first failure cannot have odd weight when \(q\equiv-1\pmod3\),
because \(q^w-1\) is then a three-adic unit for odd \(w\).

Finally choose one of the rational automorphisms of the tripod
permuting \(0,1,\infty\) and carrying \(c\) to zero, and pull back
the family E.192.2 along it.  Interior points remain interior and
the named cusp becomes the split multiplicative cusp.  The argument
then proves the assertion for every \(c\).
\end{proof}

\begin{result}{E.193}{Proved}
{A frozen ordinary reference Frobenius cannot become a boundary Frobenius at all depths}
Let \(K\) be a nonarchimedean local field with finite residue field
\(k=\mathbb F_q\) of characteristic different from three.  Fix a
reference point \(t_0\in U(K)\) whose reduction lies in \(U(k)\).
An actual point \(t\in U(K)\) reducing to a cusp cannot have its
Frobenius image geometrically conjugate to that of \(t_0\) in every
finite quotient of the geometrically pro-three arithmetic group.
The same conclusion holds for a full section locally conjugate to
the point section of \(t\).

More precisely, choose the cusp-specific family in E.192 and any
\(N\) satisfying E.192.1.  Its finite modulo-\(3^N\) representation
already distinguishes their Frobenius conjugacy classes, independently
of the nonzero ordinary valuation at the new support prime.
\end{result}

\begin{proof}
Use a tripod automorphism to place the reduction cusp at zero.
The reference elliptic fibre has good reduction and its unramified
Frobenius representation is that of the reduced smooth elliptic curve;
this follows from specialization of prime-to-residue-characteristic
torsion.  More explicitly, reduction injects the geometric
\(3^n\)-torsion into the special fibre's \(3^n\)-torsion, and both
groups have order \(3^{2n}\); the resulting isomorphism is equivariant
for the unramified Galois action.  Taking inverse limits gives the
asserted Tate-module comparison
\citep[VII.3, Proposition 3.1, p.192; VII.4, Proposition 4.1,
p.195]{AuditSilverman2009}.
Its trace is \(q+1-m\), where \(1\le m\le2q+1\).
For the boundary-disc point \(t\), the same Weierstrass model has
unit \(c_4\), positive discriminant valuation, and the split node
E.192.6.  Tate uniformization therefore gives a local exact sequence
of the form E.192.7.  Every Frobenius lift has trace \(q+1\).
These traces remain different modulo \(3^N\), exactly as in E.192.
Local geometric conjugation does not change them.

This has the following precise application to a recursive global
construction.  At a frozen support prime \(q\ne3\), one cannot
retain the Frobenius class of one fixed rational reference point
having good interior reduction while prescribing a boundary point
class in every cofinal finite three-group quotient.  Eventually the
finite detector above is among the requirements, and they conflict.
This does not invalidate E.128 or E.132, which retain only a specified
finite amount of reference data when a new prime is chosen.  Nor
does it rule out changing the higher global section data after that
prime is frozen.  It therefore does not prove nonexistence of
all-depth ordinary ghosts.

There is also an unconditional necessary representation test for any
putative full geometrically pro-three global section.  The family
E.192.2 is defined over \(\mathbb Q\), still has geometric pro-three
monodromy, and can be pulled back along that section.  At every
off-diagonal support place where A.31 identifies the section with
its actual ordinary boundary point, the corresponding cusp-specific
representation has the Tate exact sequence E.192.7 and Frobenius
trace \(q+1\).  This conclusion uses no point-matching assertion
at zero-support places or at the coefficient prime.  Such a
continuous representation may have infinite ramification; no
finite-ramification or global geometricity conclusion follows from
the detector alone.
\end{proof}

\begin{result}{E.194}{Proved}
{A two-primary detector in every characteristic different from two}
The statement of E.192 also holds with three replaced by two and
with \(\operatorname{char}(k)\ne2\).  One can take the family
\[
                   \mathcal L_t:\quad y^2=x(x+1)(x+t).
\tag{E.194.1}
\]
For any \(N\) with \(2^N>2q+1\), its cusp-specific
modulo-\(2^N\) representation separates the point from the cusp.
Its geometric image has order at most \(2^{3N-3}\).
\end{result}

\begin{proof}
All three nonzero two-torsion points
\((0,0),(-1,0),(-t,0)\) are global sections and are distinct on
\(U\).  Hence geometric monodromy is trivial modulo two; it has
determinant one and is contained in the pro-two congruence kernel
of \(\operatorname{SL}_2(\mathbb Z_2)\).
The actual arithmetic Tate-module representation therefore factors
through the geometrically pro-two fundamental group.

Here
\[
 \Delta=16t^2(t-1)^2,\qquad c_4=16(t^2-t+1).
\tag{E.194.2}
\]
At zero the node is \(y^2=x^2(x+1)\), with the two rational distinct
tangents \(y=x\) and \(y=-x\).  Thus the Tate exact sequence has
diagonal characters \(\chi_2,1\) without a quadratic twist.
Every cusp Frobenius has trace \(q+1\).  An interior fibre has trace
\(q+1-m\), where \(1\le m\le2q+1\), so reduction modulo \(2^N\)
already separates them.  The congruence kernel modulo \(2^N\) in
\(\operatorname{SL}_2\) has order \(2^{3(N-1)}\), proving the bound.
The characteristic finite nilpotent quotient can be supplied as
in E.192.10, now with \(d=3N-3\).  Tripod automorphisms handle the
other cusps, and the same proof gives the local frozen-reference
consequence of E.193 in residue characteristic different from two.
\end{proof}

The symbolic invariants, the exact residual stabilizer, and finite
point counts including characteristic two are independently checked
in \path{code/ghost_audit_elliptic_detector.py}, with its saved output
and JSON certificate.  The computations verify the displayed
algebra and finite examples; the arguments above prove the
all-finite-field statements.
